\documentclass[10pt,a4paper]{amsart}
\usepackage{amsmath,amssymb,amsthm,graphicx}
\usepackage{verbatim}
\newenvironment{reference}{\comment}{\endcomment}
%\newenvironment{reference}{}{}
\newenvironment{slogan}{\comment}{\endcomment}
\newenvironment{history}{\comment}{\endcomment}

% For commutative diagrams we use Xy-pic
\usepackage[all]{xy}

% We use 2cell for 2-commutative diagrams.
\xyoption{2cell}
% Theorem environments.
%
\theoremstyle{plain}
\newtheorem{theorem}[subsection]{Theorem}
\newtheorem{proposition}[subsection]{Proposition}
\newtheorem{lemma}[subsection]{Lemma}

\theoremstyle{definition}
\newtheorem{definition}[subsection]{Definition}
\newtheorem{example}[subsection]{Example}
\newtheorem{exercise}[subsection]{Exercise}
\newtheorem{situation}[subsection]{Situation}

\theoremstyle{remark}
\newtheorem{remark}[subsection]{Remark}
\newtheorem{remarks}[subsection]{Remarks}

\numberwithin{equation}{subsection}

% Macros
%
\def\lim{\mathop{\mathrm{lim}}\nolimits}
\def\colim{\mathop{\mathrm{colim}}\nolimits}
\def\Spec{\mathop{\mathrm{Spec}}}
\def\Hom{\mathop{\mathrm{Hom}}\nolimits}
\def\Ext{\mathop{\mathrm{Ext}}\nolimits}
\def\SheafHom{\mathop{\mathcal{H}\!\mathit{om}}\nolimits}
\def\SheafExt{\mathop{\mathcal{E}\!\mathit{xt}}\nolimits}
\def\Sch{\mathit{Sch}}
\def\Mor{\mathop{\mathrm{Mor}}\nolimits}
\def\Ob{\mathop{\mathrm{Ob}}\nolimits}
\def\Sh{\mathop{\mathit{Sh}}\nolimits}
\def\NL{\mathop{N\!L}\nolimits}
\def\CH{\mathop{\mathrm{CH}}\nolimits}
\def\proetale{{pro\text{-}\acute{e}tale}}
\def\etale{{\acute{e}tale}}
\def\QCoh{\mathit{QCoh}}
\def\Ker{\mathop{\mathrm{Ker}}}
\def\Im{\mathop{\mathrm{Im}}}
\def\Coker{\mathop{\mathrm{Coker}}}
\def\Coim{\mathop{\mathrm{Coim}}}

% Boxtimes
%
\DeclareMathSymbol{\boxtimes}{\mathbin}{AMSa}{"02}

%
% Macros for moduli stacks/spaces
%
\def\QCohstack{\mathcal{QC}\!\mathit{oh}}
\def\Cohstack{\mathcal{C}\!\mathit{oh}}
\def\Spacesstack{\mathcal{S}\!\mathit{paces}}
\def\Quotfunctor{\mathrm{Quot}}
\def\Hilbfunctor{\mathrm{Hilb}}
\def\Curvesstack{\mathcal{C}\!\mathit{urves}}
\def\Polarizedstack{\mathcal{P}\!\mathit{olarized}}
\def\Complexesstack{\mathcal{C}\!\mathit{omplexes}}
% \Pic is the operator that assigns to X its picard group, usage \Pic(X)
% \Picardstack_{X/B} denotes the Picard stack of X over B
% \Picardfunctor_{X/B} denotes the Picard functor of X over B
\def\Pic{\mathop{\mathrm{Pic}}\nolimits}
\def\Picardstack{\mathcal{P}\!\mathit{ic}}
\def\Picardfunctor{\mathrm{Pic}}
\def\Deformationcategory{\mathcal{D}\!\mathit{ef}}
\usepackage{cleveref}
\crefname{originalreference}{original reference}{original references}
\Crefname{originalreference}{Original reference}{Original references}
\makeatletter
\newlabel{sites-lemma-when-shriek}{{00XR}{source}{}{sourcepointer-1}{}}
\newlabel{sites-lemma-when-shriek@cref}{{[originalreference][0][]00XR}{[1][0][]source}{}{}{}}
\makeatother
\crefname{originalreference}{original reference}{original references}
\Crefname{originalreference}{Original reference}{Original references}
\makeatletter
\newlabel{homology-lemma-partially-defined-adjoint}{{0793}{source}{}{sourcepointer-2}{}}
\newlabel{homology-lemma-partially-defined-adjoint@cref}{{[originalreference][0][]0793}{[1][0][]source}{}{}{}}
\makeatother
\crefname{originallemma}{original lemma}{original lemmas}
\Crefname{originallemma}{Original lemma}{Original lemmas}
\makeatletter
\newlabel{lemma-module-quotient-flat}{{03EW}{source}{}{sourcepointer-3}{}}
\newlabel{lemma-module-quotient-flat@cref}{{[originallemma][0][]03EW}{[1][0][]source}{}{}{}}
\makeatother
\title{Existence of the lower shriek functor on modules for a continuous and cocontinuous functor between sites}
\author{The Stacks Project Authors}
\date{stacks-project:sites-modules.tex:local-snapshot-20261012; GNU FDL 1.2 or later}
\begin{document}
\maketitle
\noindent\textit{Editorial note.} All delivered mathematical text is the author's own, copied verbatim from the pinned Stacks Project chapter file sites-modules.tex; no statement or proof step was rewritten, reordered or removed.\par\medskip
\noindent\textit{Editorial note.} The counted span cites two external items, recorded in assembly.external\_references with canonical tags as navigation pointers.\par\medskip
\noindent\textit{Allowed standard prerequisites (editorial).}\begin{itemize}
\item The adjunction defining extension by zero and the relation g\textasciicircum{}\{-1\}G(U) = G(u(U)). Stacks Project, Modules on Sites, Lemma lemma-extension-by-zero (tag 09B5). Supplies the two identifications that make the Hom-set computation work.
\item Continuity and cocontinuity of the functor u. Stacks Project, Sites, section on continuous and cocontinuous functors. Ensures that the associated morphism of topoi g has the stated properties used in the computation.
\end{itemize}\medskip
\ifdefined\hypertarget\hypertarget{sourcepointer-1}{}\fi
\noindent\textit{Source reference.} Original reference 00XR: sources/stacks-chow/tags.txt, source lines 1229--1229. Cross-chapter standard prerequisite cited by the author inside the counted span; it is used as a pointer and its own proof is not part of this item.\par\medskip
\ifdefined\hypertarget\hypertarget{sourcepointer-2}{}\fi
\noindent\textit{Source reference.} Original reference 0793: sources/stacks-chow/tags.txt, source lines 9258--9258. Cross-chapter standard prerequisite cited by the author inside the counted span; it is used as a pointer and its own proof is not part of this item.\par\medskip
\ifdefined\hypertarget\hypertarget{sourcepointer-3}{}\fi
\noindent\textit{Source reference.} Original lemma 03EW: sources/stacks-chow/tags.txt, source lines 4471--4471. Cross-chapter standard prerequisite cited by the author inside the counted span; it is used as a pointer and its own proof is not part of this item.\par\medskip
\begin{lemma}
\label{lemma-lower-shriek-modules}
Let $u : \mathcal{C} \to \mathcal{D}$ be a continuous and cocontinuous
functor between sites. Denote
$g : \Sh(\mathcal{C}) \to \Sh(\mathcal{D})$ the associated
morphism of topoi. Let $\mathcal{O}_\mathcal{D}$ be a sheaf of rings
on $\mathcal{D}$. Set
$\mathcal{O}_\mathcal{C} = g^{-1}\mathcal{O}_\mathcal{D}$.
Hence $g$ becomes a morphism of ringed topoi with $g^* = g^{-1}$.
In this case there exists a functor
$$
g_! :
\textit{Mod}(\mathcal{O}_\mathcal{C})
\longrightarrow
\textit{Mod}(\mathcal{O}_\mathcal{D})
$$
which is left adjoint to $g^*$.
\end{lemma}

\begin{proof}
Let $U$ be an object of $\mathcal{C}$. For any
$\mathcal{O}_\mathcal{D}$-module $\mathcal{G}$ we have
\begin{align*}
\Hom_{\mathcal{O}_\mathcal{C}}(j_{U!}\mathcal{O}_U, g^{-1}\mathcal{G})
& =
g^{-1}\mathcal{G}(U) \\
& =
\mathcal{G}(u(U)) \\
& =
\Hom_{\mathcal{O}_\mathcal{D}}(j_{u(U)!}\mathcal{O}_{u(U)}, \mathcal{G})
\end{align*}
because $g^{-1}$ is described by restriction, see
Sites, Lemma \ref{sites-lemma-when-shriek}.
Of course a similar formula holds a direct sum of modules
of the form $j_{U!}\mathcal{O}_U$. By
Homology, Lemma \ref{homology-lemma-partially-defined-adjoint}
and
Lemma \ref{lemma-module-quotient-flat}
we see that $g_!$ exists.
\end{proof}

\end{document}
